% Options for packages loaded elsewhere
\PassOptionsToPackage{unicode}{hyperref}
\PassOptionsToPackage{hyphens}{url}
\PassOptionsToPackage{dvipsnames,svgnames,x11names}{xcolor}
%
\documentclass[
  11pt,
  a4paper,
]{ltjsarticle}
\usepackage{amsmath,amssymb}
\usepackage{iftex}
\ifPDFTeX
  \usepackage[T1]{fontenc}
  \usepackage[utf8]{inputenc}
  \usepackage{textcomp} % provide euro and other symbols
\else % if luatex or xetex
  \usepackage{unicode-math} % this also loads fontspec
  \defaultfontfeatures{Scale=MatchLowercase}
  \defaultfontfeatures[\rmfamily]{Ligatures=TeX,Scale=1}
\fi
\usepackage{lmodern}
\ifPDFTeX\else
  % xetex/luatex font selection
\fi
% Use upquote if available, for straight quotes in verbatim environments
\IfFileExists{upquote.sty}{\usepackage{upquote}}{}
\IfFileExists{microtype.sty}{% use microtype if available
  \usepackage[]{microtype}
  \UseMicrotypeSet[protrusion]{basicmath} % disable protrusion for tt fonts
}{}
\makeatletter
\@ifundefined{KOMAClassName}{% if non-KOMA class
  \IfFileExists{parskip.sty}{%
    \usepackage{parskip}
  }{% else
    \setlength{\parindent}{0pt}
    \setlength{\parskip}{6pt plus 2pt minus 1pt}}
}{% if KOMA class
  \KOMAoptions{parskip=half}}
\makeatother
\usepackage{xcolor}
\usepackage[margin=24mm]{geometry}
\setlength{\emergencystretch}{3em} % prevent overfull lines
\providecommand{\tightlist}{%
  \setlength{\itemsep}{0pt}\setlength{\parskip}{0pt}}
\setcounter{secnumdepth}{5}
\ifLuaTeX
  \usepackage{selnolig}  % disable illegal ligatures
\fi
\IfFileExists{bookmark.sty}{\usepackage{bookmark}}{\usepackage{hyperref}}
\IfFileExists{xurl.sty}{\usepackage{xurl}}{} % add URL line breaks if available
\urlstyle{same}
\hypersetup{
  colorlinks=true,
  linkcolor={blue},
  filecolor={Maroon},
  citecolor={Blue},
  urlcolor={blue},
  pdfcreator={LaTeX via pandoc}}

\author{}
\date{}

\begin{document}

\hypertarget{ux89e3ux7b54-04-ux5185ux7a4dux76f4ux4ea4ux5c04ux5f71}{%
\section{解答 04 ---
内積・直交・射影}\label{ux89e3ux7b54-04-ux5185ux7a4dux76f4ux4ea4ux5c04ux5f71}}

\hypertarget{section}{%
\subsection{1}\label{section}}

\[
\cos\theta=\frac{1}{\sqrt2}
\]

なので \texttt{θ=π/4} です。

\hypertarget{section-1}{%
\subsection{2}\label{section-1}}

\[
\operatorname{proj}_a b=\frac{a^Tb}{a^Ta}a=\frac52(1,1)^T.
\]

\hypertarget{section-2}{%
\subsection{3}\label{section-2}}

\texttt{q\^{}Tq=1} より

\[
P^2=qq^Tqq^T=q(q^Tq)q^T=qq^T=P.
\]

また \texttt{(qq\^{}T)\^{}T=qq\^{}T} です。射影行列の idempotent
性と対称性が確認できます。

\hypertarget{section-3}{%
\subsection{4}\label{section-3}}

\[
q_1=\frac1{\sqrt2}(1,1,0)^T.
\]

\texttt{a\_2} から \texttt{q\_1} 成分を引くと

\[
u_2=(1,0,1)^T-\frac12(1,1,0)^T=(1/2,-1/2,1)^T.
\]

そのノルムは \texttt{√(3/2)} なので正規化すれば \texttt{q\_2}
が得られます。

\hypertarget{section-4}{%
\subsection{5}\label{section-4}}

列空間上の候補点から \texttt{b}
への差に列空間方向の成分が残っていれば、その方向へ少し移動することで距離をさらに短くできます。最短点では改善可能な接線方向成分が0になるため、残差は列空間に直交します。

\end{document}
